\begin{examplebox}
\textbf{例1: 已知 $OP$ 的斜率, 求切线方程 (基础母题)}

已知椭圆 $\frac{x^2}{4}+y^2=1$, $P$ 在第一象限,
$k_{OP}=1$, 求 $P$ 处的切线方程.

\textbf{解}:
\begin{enumerate}
  \item 由 $a^2=4$, $b^2=1$, 得
  $k_lk_{OP}=-1/4$, 所以 $k_l=-1/4$.
  \item 因为 $k_{OP}=y_0/x_0=1$, 故 $y_0=x_0$.
  \item 代入椭圆:
  \[
  \frac{x_0^2}{4}+x_0^2=1
  \ \Longrightarrow\ \frac54x_0^2=1.
  \]
  由第一象限条件, 得
  $x_0=y_0=\frac2{\sqrt5}=\frac{2\sqrt5}{5}$.
  \item 用点斜式:
  \[
  y-\frac{2\sqrt5}{5}
  =-\frac14\left(x-\frac{2\sqrt5}{5}\right).
  \]
  整理得
  \[
  \boxed{x+4y-2\sqrt5=0.}
  \]
  \item 用切线公式验算:
  $\frac{x_0x}{4}+y_0y=1$,
  代入 $x_0=y_0=2/\sqrt5$ 后同样得到 $x+4y=2\sqrt5$.
\end{enumerate}
\textbf{方法提醒}: 斜率定值先求 $k_l$ 很快;
但\textbf{求切线方程还要知道经过哪一点}.

\medskip
\textbf{例2: 只给切线斜率, 求两条平行切线 (一题两解)}

椭圆 $\frac{x^2}{9}+\frac{y^2}{4}=1$, 求斜率为 $-1/2$ 的全部切线.

\textbf{解法A: 用斜率积}
\begin{enumerate}
  \item 由 $k_lk_{OP}=-4/9$、$k_l=-1/2$, 得
  \[
  k_{OP}=\frac{-4/9}{-1/2}=\frac89.
  \]
  \item 令切点 $P(x_0,y_0)$, 则 $y_0=\frac89x_0$.
  代入椭圆:
  \[
  \frac{x_0^2}{9}+\frac1{4}\left(\frac89x_0\right)^2=1
  \ \Longrightarrow\ \frac{25x_0^2}{81}=1.
  \]
  因而
  \[
  P_1=\left(\frac95,\frac85\right),\qquad
  P_2=\left(-\frac95,-\frac85\right).
  \]
  \item 由切线方程 $\frac{x_0x}{9}+\frac{y_0y}{4}=1$,
  两条切线分别为
  \[
  \boxed{x+2y=5},\qquad \boxed{x+2y=-5}.
  \]
\end{enumerate}
\textbf{解法B: 判别式速算}

设切线为 $y=-x/2+t$, 即 $x+2y=2t$.
代入椭圆并乘以 $36$:
\[
4x^2+9\left(-\frac x2+t\right)^2=36.
\]
乘以 $4$ 展开:
\[
25x^2-36tx+36t^2-144=0.
\]
相切时判别式为零:
\[
(-36t)^2-4\cdot25(36t^2-144)=0
\ \Longrightarrow\ 14400-2304t^2=0.
\]
所以 $t^2=25/4$, $t=\pm5/2$,
得到 $x+2y=\pm5$, 与上面一致.

\medskip
\textbf{例3: 从离心率直接求斜率 (快速应用)}

某椭圆长轴沿 $x$ 轴, 离心率 $e=3/5$.
若非轴切点满足 $k_{OP}=4/5$, 求切线斜率.

\textbf{解}:
\[
k_lk_{OP}=e^2-1=\frac9{25}-1=-\frac{16}{25}.
\]
因为 $k_{OP}=4/5$, 所以
\[
\boxed{k_l=\frac{-16/25}{4/5}=-\frac45.}
\]
\textbf{提醒}: 不需要先求 $a,b$ 的具体长度; 离心率已经决定轴比.

\medskip
\textbf{例4: 平行弦中点与切线的联系 (高价值变式)}

椭圆 $\frac{x^2}{9}+\frac{y^2}{4}=1$ 与直线 $y=x+1$
交于两点 $A,B$, 求弦中点 $M$ 及 $OM$ 的斜率.

\textbf{解}:
\begin{enumerate}
  \item 联立并乘以 $36$:
  \[
  4x^2+9(x+1)^2=36
  \ \Longrightarrow\ 13x^2+18x-27=0.
  \]
  \item 设两根为 $x_1,x_2$, 由韦达定理
  $x_1+x_2=-18/13$, 所以中点横坐标
  $u=(x_1+x_2)/2=-9/13$.
  \item 两端点均满足 $y=x+1$, 故中点纵坐标
  $v=u+1=4/13$. 因而
  \[
  M=\left(-\frac9{13},\frac4{13}\right).
  \]
  \item 所以
  \[
  \boxed{k_{OM}=\frac{4/13}{-9/13}=-\frac49.}
  \]
  验证: 弦斜率 $k_{AB}=1$, 所以
  $k_{AB}k_{OM}=-4/9=-b^2/a^2$.
\end{enumerate}
\textbf{变式识别}: 一组斜率为 $1$ 的平行弦, 只要中点不是中心且 $OM$ 斜率存在,
其中心连线斜率都等于 $-4/9$; 当弦退化为切线, 同一定值仍成立.

\medskip
\textbf{例5: 过切点的中心连线与切线垂直, 切点在哪里? (易错反例)}

已知非圆椭圆 $\frac{x^2}{a^2}+\frac{y^2}{b^2}=1$ ($a>b>0$),
在什么位置 $OP$ 与 $P$ 处切线垂直?

\textbf{解}:
\begin{enumerate}
  \item 若 $x_0y_0\ne0$, 两斜率存在, 且
  $k_lk_{OP}=-b^2/a^2\ne-1$. 因此\textbf{不垂直}.
  \item 若 $P=(\pm a,0)$, $OP$ 水平, 切线为 $x=\pm a$, 竖直; 二者垂直.
  \item 若 $P=(0,\pm b)$, $OP$ 竖直, 切线为 $y=\pm b$, 水平; 二者垂直.
\end{enumerate}
\[
\boxed{\text{只有四个轴顶点满足 }OP\perp l.}
\]
\textbf{重要}: 顶点处不能以 $k_lk_{OP}=-1$ 判断垂直, 因为有一条斜率不存在.

\medskip
\textbf{例6: 椭圆中心平移后, 应改用中心到切点的连线}

椭圆 $\frac{(x-1)^2}{9}+\frac{(y-2)^2}{4}=1$
的中心为 $S(1,2)$, $P(x_0,y_0)$ 为非轴切点.
求切线斜率与 $SP$ 斜率的乘积.

\textbf{解}:
\begin{enumerate}
  \item 令 $X=x-1$, $Y=y-2$; 新坐标下椭圆为 $X^2/9+Y^2/4=1$.
  \item 对应切点为 $P'(x_0-1,y_0-2)$,
  所以 $k_{SP}=\frac{y_0-2}{x_0-1}$.
  \item 切线斜率为
  $k_l=-\frac{4(x_0-1)}{9(y_0-2)}$.
  \item 相乘得到
  \[
  \boxed{k_lk_{SP}=-\frac49.}
  \]
\end{enumerate}
\textbf{提醒}: 平移中心后不能不加区分地使用原点 $O(0,0)$ 到切点的斜率.
\end{examplebox}
